% ===========================================================================
% Final thesis figure: lumped-parameter model of the dual-drive H-type gantry
% with the simulated payload absorber. Geometry adapted from
% Garcia-Herreros et al. (2013), Fig. 2.
%
% Requires in the preamble:
%   \usepackage{tikz}
%   \usetikzlibrary{arrows.meta,calc,decorations.pathmorphing}
%
% Conventions
%   * Coordinates are in drawing units; the x/y keys of the tikzpicture set the
%     unit, so the drawing scales while the labels keep the document's own font
%     and size.  At 0.3mm the drawing is 1:1 with Garcia Fig. 2 and the figure
%     is about 275pt wide, i.e. half of a two column text width (figure*).
%   * The origin is the cross arm centre of mass (the blue coordinate X).
%   * Parts on the cross arm are drawn in its rotated frame: u along the beam
%     (positive to the right), v perpendicular to it.  The measured coordinate
%     Y points LEFT in this view, so Y = -u, and the absorber, which sits at
%     +L0 in Y, is drawn to the LEFT of the payload.
%   * Changing \th (yaw) or \Yp (payload position) moves every attached part,
%     including its labels.  With \Yp > 0 the absorber ends up between the CoM
%     and the payload, and the F_Y label then needs a manual nudge.
%   * Schematic, not to scale: the absorber offset L0 and the mass sizes are
%     drawn for legibility.  The absorber rides on the payload, so it may
%     overhang the beam end.
%
% \gantrydelta: 1 adds an arrow for the absorber displacement relative to the
% payload.
% ===========================================================================
\providecommand{\gantrydelta}{1}
\begin{tikzpicture}[
    x=0.3mm, y=0.3mm,                        % original Garcia scale (1 unit
                                             % = 1pt of Garcia Fig. 2)
    line width=0.5pt, line cap=round, line join=round,
    >={Triangle[length=3.4pt,width=2.9pt]},
    dim/.style={line width=0.4pt, dash pattern=on 3pt off 1.1pt on 0.4pt off 1.1pt},
    ext/.style={line width=0.35pt},
    cl/.style={gray!55, line width=0.25pt,
               dash pattern=on 4pt off 1pt on 0.3pt off 1pt},
    ref/.style={blue!70, line width=0.3pt, dash pattern=on 1.2pt off 0.8pt},
    abs/.style={draw=black, fill=black!15, line width=0.5pt}, % absorber: shaded
    absl/.style={font=\small},                                % k_a, c_a labels
  ]
  % ================ parameters ================
  \def\th{9.25}      % yaw angle Theta [deg]
  \def\Lb{152.5}     % cross arm length
  \def\Wb{24.5}      % cross arm width
  \def\Yp{-18}       % payload position along the cross arm (u); Y = -u
  \def\Lh{26}        % payload width
  \def\Hh{24}        % payload height
  \def\rw{3.7}       % payload wheel radius
  \def\Wa{20}        % absorber width
  \def\Ha{20}        % absorber height
  \def\ga{30}        % spring and damper gap between absorber and payload
  \def\Mw{34.7}      % actuator mass box size
  \pgfmathsetmacro\yE{0.5*\Lb*sin(\th)}       % height of the beam ends
  \pgfmathsetmacro\hb{0.5*\Wb}
  \pgfmathsetmacro\vh{\hb+2*\rw+0.5*\Hh}      % payload centre height (v)
  \pgfmathsetmacro\ua{\Yp-0.5*\Lh-\ga-0.5*\Wa}  % absorber centre (u)
  \pgfmathsetmacro\vY{\vh+0.5*\Hh+30}          % height of the Y dimension line
  \pgfmathsetmacro\xM{0.5*\Lb*cos(\th)+22.3+0.5*\Mw}  % |x| of actuator centres
  % Lane for the d dimension and the F_Y arrow: on the CoM side of the payload
  % and clear of the red Y reference line, so it follows the sign of \Yp.
  \pgfmathsetmacro\ud{ifthenelse(\Yp<0,9,-9)}
  \pgfmathsetmacro\sd{ifthenelse(\Yp<0,1,-1)}

  % ===================== cross arm =====================
  \begin{scope}[rotate=\th]
    \draw (-\Lb/2,-\hb) rectangle (\Lb/2,\hb);
    \draw[cl] (-0.57*\Lb,0) -- (0.58*\Lb,0);
    \coordinate (E1) at ( \Lb/2,0);
    \coordinate (E2) at (-\Lb/2,0);

    % ==== payload m_h on its wheels ====
    \draw[fill=white] (\Yp-\Lh/2,\vh-\Hh/2) rectangle (\Yp+\Lh/2,\vh+\Hh/2);
    \draw (\Yp-7,\hb+\rw) circle (\rw)  (\Yp+7,\hb+\rw) circle (\rw);
    \node at (\Yp,\vh) {$m_h$};

    % ==== payload absorber: mass m_a, spring k_a, damper c_a ====
    \pgfmathsetmacro\xa{\ua+0.5*\Wa}          % absorber side attachment
    \pgfmathsetmacro\xh{\Yp-0.5*\Lh}          % payload side attachment
    \draw[abs] (\ua-\Wa/2,\vh-\Ha/2) rectangle (\ua+\Wa/2,\vh+\Ha/2);
    \node at (\ua,\vh) {$m_a$};
    \pgfmathsetmacro\xc{0.5*(\xa+\xh)}        % centre of the gap
    % spring k_a on top
    \draw[decorate, decoration={zigzag, segment length=4, amplitude=3.5,
          pre length=5, post length=5}] (\xa,\vh+6) -- (\xh,\vh+6);
    \node[absl, anchor=south, inner sep=1.5pt] at (\xc,\vh+10) {$k_a$};
    % damper c_a below: rod, filled piston, open cylinder, rod
    \draw (\xa,\vh-6) -- (\xc+1,\vh-6);
    \fill (\xc-1.4,\vh-8.6) rectangle (\xc+0.8,\vh-3.4);
    \draw (\xc-4,\vh-2) -- (\xc+4,\vh-2) -- (\xc+4,\vh-10) -- (\xc-4,\vh-10);
    \draw (\xc+4,\vh-6) -- (\xh,\vh-6);
    \node[absl, anchor=north, inner sep=0.5pt]        % tucked up under the
          at (\xc,\vh-10.9) {$c_a$};                   % cylinder, clear of the beam
    \ifnum\gantrydelta>0                      % displacement relative to m_h
      \draw[->] (\ua,\vh+\Ha/2+4.5) -- ++(-14,0);
      \node[anchor=south east, inner sep=1pt]
            at (\ua-3,\vh+\Ha/2+5.5) {$\delta_a$};
    \fi
    % ==== yaw reference at the payload: axis perpendicular to the beam ====
    \coordinate (P) at (\Yp,-1.5);
    \draw[ref] (P) -- ++(0,\vY+34);

    % ==== payload force F_Y ====
    \draw[->] ({\ud+\sd*11},\vh+4) -- ({\Yp+\sd*\Lh/2},\vh+4);
    \node[inner sep=2pt] at ({\ud+\sd*20},\vh+4) {$F_Y$};

    % ==== d : centreline to payload centre of mass ====
    \draw[<->] (\ud,0) -- (\ud,\vh);
    \node[inner sep=2pt] at ({\ud-\sd*4.5},0.78*\vh) {$d$};

    % ==== Y : CoM to payload (red) ====
    \begin{scope}[red]
      \draw[dim] (0,\hb+2) -- (0,\vY+6);   % the blue payload axis is the
                                            % left extension of this dimension
      \draw[->] (0,\vY) -- (\Yp,\vY);
      \node[anchor=south] at ({0.5*\Yp},\vY+1) {$Y$};
    \end{scope}

    % ==== flexible joints and L_b ====
    \draw[dim] ( \Lb/2,-\hb-2) -- ++(0,-24)   (-\Lb/2,-\hb-2) -- ++(0,-24);
    \draw[<->] (-\Lb/2,-\hb-18) -- node[pos=0.72, below=1pt] {$L_b$}
               ( \Lb/2,-\hb-18);
    \node[anchor=north west, inner sep=2.5pt] at ( \Lb/2,-\hb-19) {$k_{b1},\,c_{b1}$};
    \node[anchor=north east, inner sep=2.5pt] at (-\Lb/2,-\hb-19) {$k_{b2},\,c_{b2}$};
  \end{scope}

  % ============ Theta at the payload: vertical reference and arc ============
  \draw[ref] (P) -- ++(0,{(\vY+34)*cos(\th)});
  \draw[ref, line width=0.4pt]
        ($(P)+(90:\vY+22)$) arc[start angle=90, end angle=90+\th, radius=\vY+22];
  \node[blue, anchor=east, inner sep=2pt]
        at ($(P)+(90+\th:\vY+24)$) {$\Theta$};

  % ===================== generalized coordinates X, Theta (blue) =============
  \begin{scope}[blue]
    \draw[->] (0,-5.2) -- (0,-26);
    \node[anchor=west, inner sep=2pt] at (1.5,-21) {$X$};
    \draw (4.6,0) -- (36,0);
    \draw (32,0) arc[start angle=0, end angle=\th, radius=32];
    \node[anchor=west, inner sep=2.5pt] at (37,{36*sin(0.5*\th)}) {$\Theta$};
  \end{scope}
  % centre of mass m_b
  \fill[white] (0,0) circle (4.3);
  \fill (0,0) -- (4.3,0) arc[start angle=0, end angle=-90, radius=4.3] -- cycle
        (0,0) -- (-4.3,0) arc[start angle=180, end angle=90, radius=4.3] -- cycle;
  \draw (0,0) circle (4.3);
  \node[anchor=north east, inner sep=1pt] at (-3.4,-3.4) {$m_b$};

  % ===================== actuators, guides, forces =====================
  % \s = +1 : actuator 1 (right) ; \s = -1 : actuator 2 (left)
  \foreach \s/\i in {1/1,-1/2}{
    \pgfmathsetmacro\y{\s*\yE}                 % row height (symmetric)
    \coordinate (M\i) at (\s*\xM,\y);
    \draw (\s*\xM-\Mw/2,\y-\Mw/2) rectangle (\s*\xM+\Mw/2,\y+\Mw/2);
    \node at (M\i) {$m_\i$};

    % flexible joint: spiral spring between beam end and actuator
    \pgfmathsetmacro\xs{\s*(0.5*\Lb*cos(\th)+11)}
    \draw (\xs,\y) -- (E\i);
    \draw[line width=0.4pt] plot[domain=0:540, samples=100, smooth]  % 1.5 turns
          ({\xs+(3.4+3.9*\x/540)*cos(90*(1-\s)-\s*(\x-540))},
           {\y +(3.4+3.9*\x/540)*sin(90*(1-\s)-\s*(\x-540))});
    \draw (\xs+\s*7.3,\y) -- (\s*\xM-\s*\Mw/2,\y);

    % linear guide: two rollers against a hatched wall
    \pgfmathsetmacro\xb{\s*(\xM+\Mw/2+5.35)}
    \draw (\xb,\y+11.6) circle (5.35)  (\xb,\y-11.6) circle (5.35);
    \draw (\xb+\s*5.6,\y-20) -- (\xb+\s*5.6,\y+20);
    \foreach \k in {0,...,5}
      \draw[ext] (\xb+\s*5.6,\y+18-7.5*\k) -- ++(\s*6.5,-6.5);

    % actuator force
    \draw[->] (\s*\xM,\y+\Mw/2+30) -- (\s*\xM,\y+\Mw/2);
    \node[anchor=south] at (\s*\xM,\y+\Mw/2+31) {$F_{X\i}$};

    % measured position X_i (red)
    \pgfmathsetmacro\xr{\s*(\xM+\Mw/2+23.5)}  % clear of the wall hatching
    \draw[red] (\xr,\y+4.5) -- (\xr,\y-4.5)  (\xr,\y) -- (\xr+\s*5,\y);
    \draw[red,->] (\xr+\s*5,\y) -- (\xr+\s*5,\y-26);
    \node[red, anchor=north] at (\xr+\s*5,\y-27) {$X_\i$};
  }

  % ===================== inertial frame (top view) =====================
  \begin{scope}[shift={(72,-58)}]
    \draw[->] (0,-3.2) -- (0,-19);          % x points down
    \draw[->] (-3.2,0) -- (-19,0);          % y points left
    \draw[fill=white] (0,0) circle (3.2);   % z into the page: circle with a cross
    \draw[ext] (-2.26,-2.26) -- (2.26,2.26)  (-2.26,2.26) -- (2.26,-2.26);
    \node[anchor=east, inner sep=2pt]  at (-19,0.3) {$y$};
    \node[anchor=west, inner sep=2pt]  at (3.6,0.8) {$z$};
    \node[anchor=north, inner sep=2pt] at (0,-20)  {$x$};
  \end{scope}
\end{tikzpicture}
